\documentclass[12pt,a4paper]{article}
\usepackage[utf8]{vietnam}
\usepackage{amsmath,amssymb}
\usepackage{tabularx}
\usepackage{tikz}
\usepackage{xcolor}
\usepackage[top=0.6cm, bottom=0.6cm, left=0.6cm, right=0.6cm]{geometry}
\usepackage[dethi]{ex_test}

\SAOPTN{widthSA=2.6,heightSA=0.62}
\setlength{\parskip}{0pt}

% Dong ke cham mo chuan viet tay
\newcommand{\dongke}[1]{%
  \begingroup%
  \setlength{\parskip}{0pt}%
  \baselineskip=18.5pt%
  \vspace{1.0mm}%
  \foreach \i in {1,...,#1}{%
    \noindent{\color{gray!60}\dotfill}\par%
  }%
  \vspace{0.6mm}%
  \endgroup%
}

% Dinh nghia o bai tap tu luan trong bang
\newcommand{\oBaiTap}[3][6]{%
  \vspace{1.0mm}%
  \noindent\textbf{#2} #3\par%
  \dongke{#1}%
}

\begin{document}
\pagestyle{empty}
% =============================================================================
% ========================= MA DE 101 - MAT TRUOC (TRANG 1/2) =============
% =============================================================================
\noindent
\begin{tabularx}{\linewidth}{|X|c|}
\hline
\textbf{PHIẾU KIỂM TRA ĐỊNH KỲ TOÁN 8 - TUẦN 08 - BUỔI 02} & \textbf{Mã đề: 101 - Trang 1/2} \\
\hline
\end{tabularx}

\vspace{1.2mm}

\noindent
\begin{tabularx}{\linewidth}{|p{3.2cm}|X|}
\hline
\centering\textbf{Điểm số} & \centering\textbf{Lời nhận xét của thầy cô} \tabularnewline
\rule{0pt}{1.8cm} & \\
\hline
\end{tabularx}

\vspace{1.2mm}

\noindent\textbf{Câu 1 (0,25 điểm).} Bậc của đơn thức $A = -3x^2 y^3 z^2$ bằng bao nhiêu? \shortans{$7$}\par\vspace{0.3mm}
\noindent\textbf{Câu 2 (0,25 điểm).} Bậc của đa thức $P = 4x^3 y - 3xy^2 - 4x^3 y + 2xy^2 + 5$ là: \shortans{$3$}\par\vspace{0.3mm}
\noindent\textbf{Câu 3 (0,25 điểm).} Giá trị của $M = 3x^2 y - 2xy$ tại $x = 1, y = -3$ là: \shortans{$-3$}\par\vspace{0.3mm}
\noindent\textbf{Câu 4 (0,25 điểm).} Trong khai triển $(2x + 1)^2 = 4x^2 + ax + 1$, hệ số $a$ bằng: \shortans{$4$}\par\vspace{0.3mm}
\noindent\textbf{Câu 5 (0,25 điểm).} Tính nhanh giá trị biểu thức $101^2 - 1$ bằng: \shortans{$10\,200$}\par\vspace{0.3mm}
\noindent\textbf{Câu 6 (0,25 điểm).} Phân tích đa thức $x^2 - 6x + 9$ thành nhân tử, ta được: \shortans{$(x - 3)^2$}\par\vspace{0.3mm}
\noindent\textbf{Câu 7 (0,25 điểm).} Kết quả phép chia $(12x^3 y^2 - 8x^2 y) : (4x^2 y)$ là: \shortans{$3xy - 2$}\par\vspace{0.3mm}
\noindent\textbf{Câu 8 (0,25 điểm).} Tứ giác $ABCD$ có $\widehat{A}=75^\circ, \widehat{B}=105^\circ, \widehat{C}=80^\circ$. Số đo $\widehat{D}$ là: \shortans{$100^\circ$}\par\vspace{0.3mm}
\noindent\textbf{Câu 9 (0,25 điểm).} Hình bình hành có hai đường chéo vuông góc là hình gì? \shortans{\text{Hình thoi}}\par\vspace{0.3mm}
\noindent\textbf{Câu 10 (0,25 điểm).} Sân hình chữ nhật dài $12\text{ m}$, rộng $5\text{ m}$. Đường chéo dài: \shortans{$13\text{ m}$}\par\vspace{0.3mm}

\vspace{0.5mm}
\noindent\textbf{Câu 11 (2,0 điểm). Thực hiện phép tính (hợp lý nếu có thể):}
\vspace{0.5mm}

\noindent
\begin{tabularx}{\linewidth}{|X|X|}
\hline
\oBaiTap[10]{a)}{$2x(x^2 - 3x + 1) - x^2(2x - 4)$} & \oBaiTap[10]{b)}{$(2x + 3)^2 - (2x - 1)(2x + 1)$} \\
\hline
\oBaiTap[10]{c)}{$(24x^5 y^4 - 16x^4 y^3 + 8x^3 y^2) : 8x^3 y^2$} & \oBaiTap[10]{d)}{$(x - 3)(x + 4) - x(x + 1)$} \\
\hline
\end{tabularx}

\newpage
% =============================================================================
% ========================= MA DE 101 - MAT SAU (TRANG 2/2) ==============
% =============================================================================
\noindent
\begin{tabularx}{\linewidth}{|X|}
\hline
\centering\textbf{TOÁN 8 - TUẦN 08 - BUỔI 02 | Mã đề: 101 - Trang 2/2} \\
\hline
\end{tabularx}

\vspace{1.2mm}

\noindent\textbf{Câu 12 (2,0 điểm). Tìm $x$, biết:}
\vspace{0.5mm}

\noindent
\begin{tabularx}{\linewidth}{|X|X|}
\hline
\oBaiTap[5]{a)}{$3x(x - 4) - 3x^2 = 24$} & \oBaiTap[5]{b)}{$2x(x - 3) + 5(x - 3) = 0$} \\
\hline
\oBaiTap[5]{c)}{$(x - 1)^2 - 16 = 0$} & \oBaiTap[5]{d)}{$x^2 - 4x + 4 = 25$} \\
\hline
\end{tabularx}

\vspace{1.2mm}

\noindent
\begin{tabularx}{\linewidth}{|X|}
\hline
\vspace{0.5mm}
\textbf{Câu 13 (2,5 điểm).} Cho tam giác $ABC$ cân tại $A$, đường trung tuyến $AM$. Gọi $N$ là trung điểm của $AC$. Lấy điểm $E$ đối xứng với $M$ qua $N$.\par
\textbf{a) (1,25 điểm)} Chứng minh tứ giác $AMCE$ là hình chữ nhật.\par
\textbf{b) (1,25 điểm)} Chứng minh tứ giác $ABME$ là hình bình hành.\par
\vspace{0.4mm}
\noindent\textit{Bài giải:}\dongke{11}
\\ \hline
\end{tabularx}

\vspace{1.2mm}

\noindent
\begin{tabularx}{\linewidth}{|X|}
\hline
\vspace{0.5mm}
\textbf{Câu 14 (1,0 điểm).} Cho biểu thức $B = -x^2 + 4x + 2022$. Tìm giá trị lớn nhất của biểu thức $B$.\par
\vspace{0.4mm}
\noindent\textit{Bài giải:}\dongke{8}
\\ \hline
\end{tabularx}

\newpage
% =============================================================================
% ========================= MA DE 102 - MAT TRUOC (TRANG 1/2) =============
% =============================================================================
\noindent
\begin{tabularx}{\linewidth}{|X|c|}
\hline
\textbf{PHIẾU KIỂM TRA ĐỊNH KỲ TOÁN 8 - TUẦN 08 - BUỔI 02} & \textbf{Mã đề: 102 - Trang 1/2} \\
\hline
\end{tabularx}

\vspace{1.2mm}

\noindent
\begin{tabularx}{\linewidth}{|p{3.2cm}|X|}
\hline
\centering\textbf{Điểm số} & \centering\textbf{Lời nhận xét của thầy cô} \tabularnewline
\rule{0pt}{1.8cm} & \\
\hline
\end{tabularx}

\vspace{1.2mm}

\noindent\textbf{Câu 1 (0,25 điểm).} Bậc của đơn thức $B = 4x^3 y z^4$ bằng bao nhiêu? \shortans{$8$}\par\vspace{0.3mm}
\noindent\textbf{Câu 2 (0,25 điểm).} Bậc của đa thức $Q = 5x^2 y^2 + 2x^3 - 5x^2 y^2 - x^3 + 8$ là: \shortans{$3$}\par\vspace{0.3mm}
\noindent\textbf{Câu 3 (0,25 điểm).} Giá trị của $N = 2x^2 y - 4xy$ tại $x = 1, y = -2$ là: \shortans{$4$}\par\vspace{0.3mm}
\noindent\textbf{Câu 4 (0,25 điểm).} Trong khai triển $(3x - 1)^2 = 9x^2 - bx + 1$, hệ số $b$ bằng: \shortans{$6$}\par\vspace{0.3mm}
\noindent\textbf{Câu 5 (0,25 điểm).} Tính nhanh giá trị biểu thức $201^2 - 1$ bằng: \shortans{$40\,400$}\par\vspace{0.3mm}
\noindent\textbf{Câu 6 (0,25 điểm).} Phân tích đa thức $x^2 - 8x + 16$ thành nhân tử, ta được: \shortans{$(x - 4)^2$}\par\vspace{0.3mm}
\noindent\textbf{Câu 7 (0,25 điểm).} Kết quả phép chia $(15x^3 y^2 - 10x^2 y) : (5x^2 y)$ là: \shortans{$3xy - 2$}\par\vspace{0.3mm}
\noindent\textbf{Câu 8 (0,25 điểm).} Tứ giác $ABCD$ có $\widehat{A}=65^\circ, \widehat{B}=115^\circ, \widehat{C}=70^\circ$. Số đo $\widehat{D}$ là: \shortans{$110^\circ$}\par\vspace{0.3mm}
\noindent\textbf{Câu 9 (0,25 điểm).} Hình bình hành có hai đường chéo bằng nhau là hình gì? \shortans{\text{Hình chữ nhật}}\par\vspace{0.3mm}
\noindent\textbf{Câu 10 (0,25 điểm).} Sân hình chữ nhật dài $8\text{ m}$, rộng $6\text{ m}$. Đường chéo dài: \shortans{$10\text{ m}$}\par\vspace{0.3mm}

\vspace{0.5mm}
\noindent\textbf{Câu 11 (2,0 điểm). Thực hiện phép tính (hợp lý nếu có thể):}
\vspace{0.5mm}

\noindent
\begin{tabularx}{\linewidth}{|X|X|}
\hline
\oBaiTap[10]{a)}{$3x(x^2 - 2x + 1) - x^2(3x - 5)$} & \oBaiTap[10]{b)}{$(2x + 1)^2 - (2x - 3)(2x + 3)$} \\
\hline
\oBaiTap[10]{c)}{$(20x^5 y^4 - 15x^4 y^3 + 5x^3 y^2) : 5x^3 y^2$} & \oBaiTap[10]{d)}{$(x - 2)(x + 5) - x(x + 3)$} \\
\hline
\end{tabularx}

\newpage
% =============================================================================
% ========================= MA DE 102 - MAT SAU (TRANG 2/2) ==============
% =============================================================================
\noindent
\begin{tabularx}{\linewidth}{|X|}
\hline
\centering\textbf{TOÁN 8 - TUẦN 08 - BUỔI 02 | Mã đề: 102 - Trang 2/2} \\
\hline
\end{tabularx}

\vspace{1.2mm}

\noindent\textbf{Câu 12 (2,0 điểm). Tìm $x$, biết:}
\vspace{0.5mm}

\noindent
\begin{tabularx}{\linewidth}{|X|X|}
\hline
\oBaiTap[5]{a)}{$4x(x - 3) - 4x^2 = 24$} & \oBaiTap[5]{b)}{$3x(x - 2) + 4(x - 2) = 0$} \\
\hline
\oBaiTap[5]{c)}{$(x - 2)^2 - 25 = 0$} & \oBaiTap[5]{d)}{$x^2 - 6x + 9 = 36$} \\
\hline
\end{tabularx}

\vspace{1.2mm}

\noindent
\begin{tabularx}{\linewidth}{|X|}
\hline
\vspace{0.5mm}
\textbf{Câu 13 (2,5 điểm).} Cho tam giác $ABC$ cân tại $A$, đường cao $AH$. Gọi $I$ là trung điểm của $AC$. Lấy điểm $K$ đối xứng với $H$ qua $I$.\par
\textbf{a) (1,25 điểm)} Chứng minh tứ giác $AHCK$ là hình chữ nhật.\par
\textbf{b) (1,25 điểm)} Chứng minh tứ giác $ABHK$ là hình bình hành.\par
\vspace{0.4mm}
\noindent\textit{Bài giải:}\dongke{11}
\\ \hline
\end{tabularx}

\vspace{1.2mm}

\noindent
\begin{tabularx}{\linewidth}{|X|}
\hline
\vspace{0.5mm}
\textbf{Câu 14 (1,0 điểm).} Cho biểu thức $B = -x^2 + 6x + 2017$. Tìm giá trị lớn nhất của biểu thức $B$.\par
\vspace{0.4mm}
\noindent\textit{Bài giải:}\dongke{8}
\\ \hline
\end{tabularx}

\newpage
% =============================================================================
% ========================= MA DE 103 - MAT TRUOC (TRANG 1/2) =============
% =============================================================================
\noindent
\begin{tabularx}{\linewidth}{|X|c|}
\hline
\textbf{PHIẾU KIỂM TRA ĐỊNH KỲ TOÁN 8 - TUẦN 08 - BUỔI 02} & \textbf{Mã đề: 103 - Trang 1/2} \\
\hline
\end{tabularx}

\vspace{1.2mm}

\noindent
\begin{tabularx}{\linewidth}{|p{3.2cm}|X|}
\hline
\centering\textbf{Điểm số} & \centering\textbf{Lời nhận xét của thầy cô} \tabularnewline
\rule{0pt}{1.8cm} & \\
\hline
\end{tabularx}

\vspace{1.2mm}

\noindent\textbf{Câu 1 (0,25 điểm).} Bậc của đơn thức $C = -5x^2 y^4 z$ bằng bao nhiêu? \shortans{$7$}\par\vspace{0.3mm}
\noindent\textbf{Câu 2 (0,25 điểm).} Bậc của đa thức $R = 3x^4 y - 2x^2 y + 1 - 3x^4 y + 5x^2 y$ là: \shortans{$3$}\par\vspace{0.3mm}
\noindent\textbf{Câu 3 (0,25 điểm).} Giá trị của $P = 4x^2 y - 3xy$ tại $x = -1, y = 2$ là: \shortans{$14$}\par\vspace{0.3mm}
\noindent\textbf{Câu 4 (0,25 điểm).} Trong khai triển $(2x + 3)^2 = 4x^2 + cx + 9$, hệ số $c$ bằng: \shortans{$12$}\par\vspace{0.3mm}
\noindent\textbf{Câu 5 (0,25 điểm).} Tính nhanh giá trị biểu thức $102^2 - 4$ bằng: \shortans{$10\,400$}\par\vspace{0.3mm}
\noindent\textbf{Câu 6 (0,25 điểm).} Phân tích đa thức $x^2 + 10x + 25$ thành nhân tử, ta được: \shortans{$(x + 5)^2$}\par\vspace{0.3mm}
\noindent\textbf{Câu 7 (0,25 điểm).} Kết quả phép chia $(18x^4 y^2 - 12x^3 y) : (6x^3 y)$ là: \shortans{$3xy - 2$}\par\vspace{0.3mm}
\noindent\textbf{Câu 8 (0,25 điểm).} Tứ giác $ABCD$ có $\widehat{A}=80^\circ, \widehat{B}=100^\circ, \widehat{C}=95^\circ$. Số đo $\widehat{D}$ là: \shortans{$85^\circ$}\par\vspace{0.3mm}
\noindent\textbf{Câu 9 (0,25 điểm).} Hình chữ nhật có hai cạnh kề bằng nhau là hình gì? \shortans{\text{Hình vuông}}\par\vspace{0.3mm}
\noindent\textbf{Câu 10 (0,25 điểm).} Sân hình chữ nhật dài $15\text{ m}$, rộng $8\text{ m}$. Đường chéo dài: \shortans{$17\text{ m}$}\par\vspace{0.3mm}

\vspace{0.5mm}
\noindent\textbf{Câu 11 (2,0 điểm). Thực hiện phép tính (hợp lý nếu có thể):}
\vspace{0.5mm}

\noindent
\begin{tabularx}{\linewidth}{|X|X|}
\hline
\oBaiTap[10]{a)}{$2x(x^2 - 4x + 2) - x^2(2x - 6)$} & \oBaiTap[10]{b)}{$(2x + 5)^2 - (2x - 2)(2x + 2)$} \\
\hline
\oBaiTap[10]{c)}{$(18x^5 y^4 - 12x^4 y^3 + 6x^3 y^2) : 6x^3 y^2$} & \oBaiTap[10]{d)}{$(x - 4)(x + 5) - x(x + 1)$} \\
\hline
\end{tabularx}

\newpage
% =============================================================================
% ========================= MA DE 103 - MAT SAU (TRANG 2/2) ==============
% =============================================================================
\noindent
\begin{tabularx}{\linewidth}{|X|}
\hline
\centering\textbf{TOÁN 8 - TUẦN 08 - BUỔI 02 | Mã đề: 103 - Trang 2/2} \\
\hline
\end{tabularx}

\vspace{1.2mm}

\noindent\textbf{Câu 12 (2,0 điểm). Tìm $x$, biết:}
\vspace{0.5mm}

\noindent
\begin{tabularx}{\linewidth}{|X|X|}
\hline
\oBaiTap[5]{a)}{$2x(x - 5) - 2x^2 = 20$} & \oBaiTap[5]{b)}{$4x(x - 4) + 3(x - 4) = 0$} \\
\hline
\oBaiTap[5]{c)}{$(x - 3)^2 - 36 = 0$} & \oBaiTap[5]{d)}{$x^2 - 8x + 16 = 49$} \\
\hline
\end{tabularx}

\vspace{1.2mm}

\noindent
\begin{tabularx}{\linewidth}{|X|}
\hline
\vspace{0.5mm}
\textbf{Câu 13 (2,5 điểm).} Cho tam giác $MNP$ cân tại $M$, đường trung tuyến $MI$. Gọi $K$ là trung điểm của $MP$. Lấy điểm $E$ đối xứng với $I$ qua $K$.\par
\textbf{a) (1,25 điểm)} Chứng minh tứ giác $MIPE$ là hình chữ nhật.\par
\textbf{b) (1,25 điểm)} Chứng minh tứ giác $MNIE$ là hình bình hành.\par
\vspace{0.4mm}
\noindent\textit{Bài giải:}\dongke{11}
\\ \hline
\end{tabularx}

\vspace{1.2mm}

\noindent
\begin{tabularx}{\linewidth}{|X|}
\hline
\vspace{0.5mm}
\textbf{Câu 14 (1,0 điểm).} Cho biểu thức $B = -x^2 + 8x + 2010$. Tìm giá trị lớn nhất của biểu thức $B$.\par
\vspace{0.4mm}
\noindent\textit{Bài giải:}\dongke{8}
\\ \hline
\end{tabularx}

\newpage
% =============================================================================
% ========================= MA DE 104 - MAT TRUOC (TRANG 1/2) =============
% =============================================================================
\noindent
\begin{tabularx}{\linewidth}{|X|c|}
\hline
\textbf{PHIẾU KIỂM TRA ĐỊNH KỲ TOÁN 8 - TUẦN 08 - BUỔI 02} & \textbf{Mã đề: 104 - Trang 1/2} \\
\hline
\end{tabularx}

\vspace{1.2mm}

\noindent
\begin{tabularx}{\linewidth}{|p{3.2cm}|X|}
\hline
\centering\textbf{Điểm số} & \centering\textbf{Lời nhận xét của thầy cô} \tabularnewline
\rule{0pt}{1.8cm} & \\
\hline
\end{tabularx}

\vspace{1.2mm}

\noindent\textbf{Câu 1 (0,25 điểm).} Bậc của đơn thức $D = 7x^3 y^2 z^3$ bằng bao nhiêu? \shortans{$8$}\par\vspace{0.3mm}
\noindent\textbf{Câu 2 (0,25 điểm).} Bậc của đa thức $S = 6x y^3 - 4x^2 + 9 - 6x y^3 + x^2$ là: \shortans{$2$}\par\vspace{0.3mm}
\noindent\textbf{Câu 3 (0,25 điểm).} Giá trị của $Q = 2x^2 y - 5xy$ tại $x = -1, y = 3$ là: \shortans{$21$}\par\vspace{0.3mm}
\noindent\textbf{Câu 4 (0,25 điểm).} Trong khai triển $(3x - 2)^2 = 9x^2 - dx + 4$, hệ số $d$ bằng: \shortans{$12$}\par\vspace{0.3mm}
\noindent\textbf{Câu 5 (0,25 điểm).} Tính nhanh giá trị biểu thức $99^2 - 1$ bằng: \shortans{$9\,800$}\par\vspace{0.3mm}
\noindent\textbf{Câu 6 (0,25 điểm).} Phân tích đa thức $x^2 + 12x + 36$ thành nhân tử, ta được: \shortans{$(x + 6)^2$}\par\vspace{0.3mm}
\noindent\textbf{Câu 7 (0,25 điểm).} Kết quả phép chia $(14x^3 y^3 - 7x^2 y^2) : (7x^2 y^2)$ là: \shortans{$2xy - 1$}\par\vspace{0.3mm}
\noindent\textbf{Câu 8 (0,25 điểm).} Tứ giác $ABCD$ có $\widehat{A}=70^\circ, \widehat{B}=110^\circ, \widehat{C}=90^\circ$. Số đo $\widehat{D}$ là: \shortans{$90^\circ$}\par\vspace{0.3mm}
\noindent\textbf{Câu 9 (0,25 điểm).} Hình thoi có một góc vuông là hình gì? \shortans{\text{Hình vuông}}\par\vspace{0.3mm}
\noindent\textbf{Câu 10 (0,25 điểm).} Sân hình chữ nhật dài $20\text{ m}$, rộng $15\text{ m}$. Đường chéo dài: \shortans{$25\text{ m}$}\par\vspace{0.3mm}

\vspace{0.5mm}
\noindent\textbf{Câu 11 (2,0 điểm). Thực hiện phép tính (hợp lý nếu có thể):}
\vspace{0.5mm}

\noindent
\begin{tabularx}{\linewidth}{|X|X|}
\hline
\oBaiTap[10]{a)}{$3x(x^2 - 3x + 2) - x^2(3x - 8)$} & \oBaiTap[10]{b)}{$(2x + 2)^2 - (2x - 1)(2x + 1)$} \\
\hline
\oBaiTap[10]{c)}{$(28x^5 y^4 - 21x^4 y^3 + 7x^3 y^2) : 7x^3 y^2$} & \oBaiTap[10]{d)}{$(x - 5)(x + 7) - x(x + 2)$} \\
\hline
\end{tabularx}

\newpage
% =============================================================================
% ========================= MA DE 104 - MAT SAU (TRANG 2/2) ==============
% =============================================================================
\noindent
\begin{tabularx}{\linewidth}{|X|}
\hline
\centering\textbf{TOÁN 8 - TUẦN 08 - BUỔI 02 | Mã đề: 104 - Trang 2/2} \\
\hline
\end{tabularx}

\vspace{1.2mm}

\noindent\textbf{Câu 12 (2,0 điểm). Tìm $x$, biết:}
\vspace{0.5mm}

\noindent
\begin{tabularx}{\linewidth}{|X|X|}
\hline
\oBaiTap[5]{a)}{$5x(x - 3) - 5x^2 = 30$} & \oBaiTap[5]{b)}{$5x(x - 1) + 2(x - 1) = 0$} \\
\hline
\oBaiTap[5]{c)}{$(x - 4)^2 - 49 = 0$} & \oBaiTap[5]{d)}{$x^2 - 10x + 25 = 64$} \\
\hline
\end{tabularx}

\vspace{1.2mm}

\noindent
\begin{tabularx}{\linewidth}{|X|}
\hline
\vspace{0.5mm}
\textbf{Câu 13 (2,5 điểm).} Cho tam giác $DEF$ cân tại $D$, đường cao $DH$. Gọi $M$ là trung điểm của $DF$. Lấy điểm $K$ đối xứng với $H$ qua $M$.\par
\textbf{a) (1,25 điểm)} Chứng minh tứ giác $DHKF$ là hình chữ nhật.\par
\textbf{b) (1,25 điểm)} Chứng minh tứ giác $DEHK$ là hình bình hành.\par
\vspace{0.4mm}
\noindent\textit{Bài giải:}\dongke{11}
\\ \hline
\end{tabularx}

\vspace{1.2mm}

\noindent
\begin{tabularx}{\linewidth}{|X|}
\hline
\vspace{0.5mm}
\textbf{Câu 14 (1,0 điểm).} Cho biểu thức $B = -x^2 + 10x + 2001$. Tìm giá trị lớn nhất của biểu thức $B$.\par
\vspace{0.4mm}
\noindent\textit{Bài giải:}\dongke{8}
\\ \hline
\end{tabularx}

\newpage
\end{document}
